\documentclass[10pt,a4paper]{amsart}
\usepackage[margin=19mm]{geometry}
\usepackage{amsmath,amssymb,mathtools}
\usepackage[scr=rsfs,cal=euler]{mathalfa}
\usepackage{xfrac}
\usepackage[unicode,hidelinks,bookmarks=false]{hyperref}
\newcommand{\ie}{i.e.\ }
\newcommand{\cf}{cf.\ }
\newcommand{\resp}{respectively\ }
\newcommand{\Subsection}{\subsection}
\providecommand{\nobreakdash}{\nobreak\hbox{-}}
\setlength{\emergencystretch}{2em}
\multlinegap0pt
\usepackage{bm}
\usepackage{bbm}
\usepackage{dsfont}
\usepackage{xspace}
\usepackage{cancel}
\usepackage{mathtools}
\usepackage{amsthm}
\makeatletter
\@ifundefined{corollary}{\newtheorem{corollary}{Corollary}}{}
\@ifundefined{lemma}{\newtheorem{lemma}[corollary]{Lemma}}{}
\@ifundefined{remark}{\newtheorem{remark}[corollary]{Remark}}{}
\makeatother
\def\B {\mathbb B}
\def\C {\mathbb C}
\def\EE {\mathrm{E}}
\def\Lin {\mathop {\rm Lin}\nolimits}
\def\M {\mathbb M^{n\times n}_{\rm sym}}
\def\R {\mathbb R}
\def\de {\mathrm{d}}
\title[A rate-independent model for geomaterials under compression ]{A rate-independent model for geomaterials under compression coupling strain gradient plasticity and damage}
\author{Maicol Caponi, Vito Crismale}
\date{Source: arXiv:2311.13863v1 (CC BY 4.0)}
\begin{document}
\maketitle

\noindent\emph{Source and licence.} A rate-independent model for geomaterials under compression coupling strain gradient plasticity and damage. Version arXiv:2311.13863v1 (CC BY 4.0), distributed under the Creative Commons Attribution 4.0 International licence, as recorded in the arXiv licence metadata for this exact version and shown on the abstract page https://arxiv.org/abs/2311.13863v1. Full rights notice: licenses/arxiv-2311.13863-CC-BY-4.0.txt.\par\smallskip

\noindent\textit{Editorial scope.} Counted unit: the Remark whose source label is \texttt{None} in the article (source0.tex lines 1002--1004, source label \texttt{None}), delivered together with the article's own complete original proof of it (source0.tex lines 1006--1106, one \texttt{proof} environment of 101 lines). No mathematics is written, repaired or re-derived: statement and proof are the article's own text line for line. Required context retained uncounted: eq:w (lines 225--227); eq:d (lines 241--243); eq:C2 (lines 250--253); eq:B1 (lines 263--267); eq:B2 (lines 263--267); eq:B3 (lines 263--267); rem:BCS (lines 269--273); eq:K (lines 290--292); eq:rH (lines 306--308); lem:Cdelta (lines 363--368); coro:H1norm (lines 398--403); lem:measurability (lines 513--519); eq:unif\_bound (lines 515--517); rem:unif\_bound (lines 589--600); eq:st\_1 (lines 867--869); eq:st\_3 (lines 875--878); lem:gronw (lines 965--974); eq:cont1 (lines 995--999); eq:cont2 (lines 995--999). Every definition, display and label of those lines is retained unchanged. Omitted from this package and not redistributed: 1--224, 228--240, 244--249, 254--262, 268--268, 274--289, 293--305, 309--362, 369--397, 404--512, 518--588, 601--866, 870--874, 879--964, 975--994, 1000--1001, 1005--1005, 1107--2395. Nothing inside a retained excerpt is omitted or altered: every retained body is delivered exactly as the article's lines are written, so no omission receipt of any kind accompanies this package. Citations: the retained excerpts carry no citation command at all, so no bibliography entry of the article is transcribed and no reference is redistributed. Reference adaptation: none was needed and none was made; every reference inside the delivered unit resolves inside it, so no unresolved reference or citation is produced. Printed slips kept literal: no printed word, symbol, hypothesis or number of the counted statement or of its proof was altered. Layout and class: the article's own class is \texttt{amsart} with packages including inputenc, amscd, amsfonts, amsmath, amssymb, amsthm, bbm, dsfont, centernot, mathtools, color, mathrsfs, stackrel, geometry; the wrapper is the lane's pinned \texttt{amsart} editorial wrapper at 10pt on A4, which restates the theorem-like environments the retained text needs and adds the article's own macro definitions that the retained text needs (\texttt{\textbackslash B}, \texttt{\textbackslash C}, \texttt{\textbackslash EE}, \texttt{\textbackslash Lin}, \texttt{\textbackslash M}, \texttt{\textbackslash R}, \texttt{\textbackslash de}), with extra wrapper packages (bm, bbm, dsfont, xspace, cancel, mathtools). The delivered numbering is the unit's own. Assets: no retained span contains a figure environment, a graphics-inclusion command, a PDF-page inclusion, a table or a TikZ picture; no image file is copied into this package, the package declares no asset dependency and no image file is redistributed with it. Licence: version arXiv:2311.13863v1 of this article is distributed under the Creative Commons Attribution 4.0 International licence, as recorded in the arXiv licence metadata for this exact version and shown on the abstract page https://arxiv.org/abs/2311.13863v1. A full rights notice is delivered at licenses/arxiv-2311.13863-CC-BY-4.0.txt. Characters: every retained excerpt is pure ASCII, so the delivered engine, XeLaTeX with its default Latin Modern text font, reports no missing character.
\begin{equation}\label{eq:w}
w\in H^1((0,T);H^1(\Omega;\R^n)).
\end{equation}

\begin{align}
&d\in C^{1,1}([0,1];[0,\infty)),\qquad \dot d(0)\le 0\label{eq:d}
\end{align}

\begin{align}
&\C_{ijkl}=\C_{klij}=\C_{jikl}\quad\text{for all $i,j,k,l=1,\dots,n$},\label{eq:C1}\\
&\gamma_1|\xi|^2\le \C\xi:\xi\le\gamma_2|\xi|^2\quad\text{for all $\xi\in\M$},\label{eq:C2}
\end{align}

\begin{align}
&\B\in C^{1,1}([0,1];\Lin(\M;\M)),\qquad \dot{\B}(0)\xi:\xi\le 0\quad\text{for all $\xi\in\M$},\label{eq:B1}\\
&\B_{ijkl}(\alpha)=\B_{klij}(\alpha)=\B_{jikl}(\alpha)\quad\text{for all $i,j,k,l=1,\dots,n$ and $\alpha\in [0,1]$},\label{eq:B2}\\
&\B(\alpha)\xi:\xi\ge 0\quad\text{for all $\alpha\in[0,1]$ and $\xi\in\M$}.\label{eq:B3}
\end{align}

\begin{remark}\label{rem:BCS}
Thanks to~\eqref{eq:B2} and~\eqref{eq:B3}, we derive the Cauchy-Schwartz Inequality \begin{equation*}
|\B(\alpha)\xi_1:\xi_2|^2\le (\B(\alpha)\xi_1:\xi_1)(\B(\alpha)\xi_2:\xi_2)\quad\text{for all $\alpha\in[0,1]$ and $\xi_1,\xi_2\in\M$}.
\end{equation*}
\end{remark}

\begin{equation}\label{eq:K}
\{\sigma\in\M\,:\,|\sigma|\le r_H\}\subset K
\end{equation}

\begin{equation}\label{eq:rH}
r_H|\xi|\le H(\xi)\quad\text{for all $\xi\in\M$}.
\end{equation}

\begin{lemma}\label{lem:Cdelta}
For all $\theta\in[2,6)$ and $\delta>0$ there exists a constant $C_{\theta,\delta}>0$ such that
\begin{equation*}
\|u\|_\theta^2\le C_{\theta,\delta}\|u\|_1^2+\delta\|\nabla u\|_2^2\quad\text{for all $u\in H^1(\Omega;\R^d)$}.
\end{equation*}
\end{lemma}

\begin{corollary}\label{coro:H1norm}
There exists a constant $C>0$ such that
\begin{equation}\label{eq:H1norm}
\|u\|_{H^1}^2\le C\left(\|u\|_1^2+\|\nabla u\|_2^2\right)\quad\text{for all $u\in H^1(\Omega;\R^d)$}.
\end{equation}
\end{corollary}

\begin{lemma}\label{lem:measurability}
Assume~\eqref{eq:w}--\eqref{eq:B3} and~\eqref{eq:K}. Let $(\alpha,u,e,p)$ from $[0,T]$ into $H^1(\Omega;[0,1])\times\mathcal A$ be a quadruple satisfying {\em (qs0)} and {\em (qs1)}, and such that $p\colon [0,T]\to H^1(\Omega;\M)$ has bounded $\mathcal H$-variation. Then there exists a constant $C>0$, independent on $t\in[0,T]$, such that
\begin{equation}\label{eq:unif_bound}
\|\alpha(t)\|_{H^1}+\|u(t)\|_{H^1}+\|e(t)\|_2+\|p(t)\|_{H^1}\le C\quad\text{for all $t\in[0,T]$}.
\end{equation}
Moreover, there exists a countable set $N\subset [0,T]$ such that the quadruple $(\alpha,u,e,p)$ is weakly continuous from $[0,T]\setminus N$ into $H^1(\Omega)\times \mathcal A$. In particular, the quadruple $(\alpha,u,e,p)$ is strongly measurable from $[0,T]$ into $H^1(\Omega)\times \mathcal A$.
\end{lemma}

\begin{remark}\label{rem:unif_bound}
If $(\alpha,u,e,p)$ from $[0,T]$ into $H^1(\Omega;[0,1])\times \mathcal A$ is a global quasistatic evolution, then we can improve the estimate in~\eqref{eq:unif_bound}. Indeed, by Lemma~\ref{lem:measurability} we have that the function $t\mapsto \|e(t)\|_2$ is bounded and measurable on $[0,T]$. In particular, if we define 
$$M\coloneqq\sup_{t\in[0,T]}\|e(t)\|_2<\infty,$$
by (qs2) we derive
$$\frac{\gamma_1}{2}M^2\le \mathcal E(\alpha(0),e(0),p(0))+\gamma_2 M\int_0^T\|\EE \dot w(s)\|_2\,\de s.$$
This implies the existence of a constant $C_1$, which depends only on $w$, $d$, $\C$, $\B$, and on the initial datum $(\alpha(0),e(0),p(0))$, satisfying
$$M=\sup_{t\in[0,T]}\|e(t)\|_2\le C_1.$$
Therefore, by using again (qs2) we derive the existence of a further constant $C_2$, depending only on $w$, $d$, $\C$, $\B$, $K$, and the initial datum $(\alpha(0),u(0),e(0),p(0))$, such that
\begin{equation*}
\|\alpha(t)\|_{H^1}+\|u(t)\|_{H^1}+\|e(t)\|_2+\|p(t)\|_{H^1}\le C_2\quad\text{for all $t\in[0,T]$}.
\end{equation*}     
\end{remark}

\begin{equation}\label{eq:st_1}
\partial_\alpha \mathcal E(\alpha(t),e(t),p(t))[\beta]\ge 0\quad\text{for all $\beta\in H^1(\Omega)$ with $\beta\le 0$}.
\end{equation}

\begin{align}
&\mathcal Q(e(t))+\mathcal Q(\eta-e(t))+\tilde Q(\alpha(t),p(t))+\tilde Q(\alpha(t),q-p(t))+\|\nabla p(t)\|_2^2+\|\nabla q-\nabla p(t)\|_2^2\nonumber\\
&\le \mathcal Q(\eta)+\tilde Q(\alpha(t),q)+\|\nabla q\|_2^2+\mathcal H(q-p(t)).\label{eq:st_3}
\end{align}

\begin{lemma}\label{lem:gronw}
Let $a,b\in\R$ with $a<b$. Let $f\colon [a,b]\to [0,\infty)$ be a bounded measurable function, let $g\colon [a,b]\to [0,\infty)$ be a non decreasing function, and let $h\colon [a,b]\to [0,\infty)$ be an integrable function. Suppose that
\begin{align}\label{eq:gronw1}
f(t)^2\le g(t)^2+\int_a^tf(r)h(r)\,\de r+\left(\int_a^t h(r)\,\de r\right)^2\quad\text{for all $t\in[a,b]$}.
\end{align}
Then, 
\begin{align}\label{eq:gronw2}
f(t)\le g(t)+\frac{3}{2}\int_a^th(r)\,\de r\quad\text{for all $t\in[a,b]$}.    
\end{align}
\end{lemma}

\begin{align}
\|\alpha(t_2)-\alpha(t_1)\|_{H^1}&+\|e(t_2)-e(t_1)\|_2+\|p(t_2)-p(t_1)\|_{H^1}\nonumber\\
&\le C\left(\|\alpha(t_2)-\alpha(t_1)\|_1+\int_{t_1}^{t_2}\|\EE \dot w(r)\|_2\,\de r\right),\label{eq:cont1}\\
\|u(t_2)-u(t_1)\|_{H^1}&\le C\left(\|\alpha(t_2)-\alpha(t_1)\|_1+\int_{t_1}^{t_2}\|\dot w(r)\|_{H^1}\,\de r\right).\label{eq:cont2}
\end{align}

\begin{remark}
Therefore, every global quasistatic evolution $(\alpha,u,e,p)$ is (strongly) continuous from $[0,T]$ into $H^1(\Omega)\times\mathcal A$, except for a countable subset of $[0,T]$, which is the set of discontinuity points of $\alpha$ with respect to the convergence in $L^1(\Omega)$. Moreover, if $\alpha\in AC([0,T];L^1(\Omega))$, then $(\alpha,u,e,p)$ are absolutely continuous from $[0,T]$ into $H^1(\Omega)\times \mathcal A$.
\end{remark}

\begin{proof}
We fix $t_1,t_2\in[0,T]$ with $t_1<t_2$. We test~\eqref{eq:st_1} in $t_1$ with $\beta\coloneqq\alpha(t_2)-\alpha(t_1)$ and we use the identity
$$\|\nabla \alpha(t_2)\|_2^2-\|\nabla \alpha(t_2)-\nabla\alpha(t_1)\|_2^2-\|\nabla \alpha(t_1)\|_2^2=2(\nabla\alpha(t_1),\nabla\alpha(t_2)-\nabla\alpha(t_1))_2,$$
to derive that
\begin{align}
&\|\nabla \alpha(t_2)-\nabla\alpha(t_1)\|_2^2\nonumber\\
&\le \|\nabla \alpha(t_2)\|_2^2-\|\nabla \alpha(t_1)\|_2^2+(\dot{d}(\alpha(t_1)),\alpha(t_2)-\alpha(t_1))_2+(\dot{\B}(\alpha(t_1))(\alpha(t_2)-\alpha(t_1))p(t_1),p(t_1))_2.\label{eq:lowbou1}  
\end{align}
We now test~\eqref{eq:st_3} in $t_1$ with
\begin{equation*}
(v,\eta,q)\coloneqq(u(t_2)-w(t_2)+w(t_1),e(t_2)-\EE w(t_2)+\EE w(t_1),p(t_2))\in \mathcal A(w(t_1)),
\end{equation*}
and we get
\begin{align*}
&\mathcal Q(e(t_1))+\mathcal Q(e(t_2)-e(t_1)-\EE(w(t_2)-w(t_1)))\\
&\quad+\tilde Q(\alpha(t_1),p(t_1))+\tilde Q(\alpha(t_1),p(t_2)-p(t_1))+\|\nabla p(t_1)\|_2^2+\|\nabla p(t_2)-\nabla p(t_1)\|_2^2\\
&\le \mathcal Q(e(t_2)-\EE(w(t_2)-w(t_1)))+\tilde Q(\alpha(t_1),p(t_2))+\|\nabla p(t_2)\|_2^2+\mathcal H(p(t_2)-p(t_1)).
\end{align*}
In particular
\begin{align}
&\mathcal Q(e(t_2)-e(t_1))+\tilde Q(\alpha(t_1),p(t_2)-p(t_1))+\|\nabla p(t_2)-\nabla p(t_1)\|_2^2\nonumber\\
&\le \mathcal Q(e(t_2))-\mathcal Q(e(t_1))-(\C e(t_1),\EE(w(t_2)-w(t_1)))_2\nonumber\\
&\quad+\tilde Q(\alpha(t_1),p(t_2))-\tilde Q(\alpha(t_1),p(t_1))+\|\nabla p(t_2)\|_2^2-\|\nabla p(t_1)\|_2^2+\mathcal H(p(t_2)-p(t_1)).\label{eq:lowbou2}
\end{align}
Moreover, by (qs2) evaluated in $t_1$ and $t_2$, we have that 
\begin{align}
\mathcal H(p(t_2)-p(t_1))&\le \mathcal V_{\mathcal H}(p;t_1,t_2)\nonumber\\
&= \mathcal E(\alpha(t_1),e(t_1),p(t_1))-\mathcal E(\alpha(t_2),e(t_2),p(t_2))+\int_{t_1}^{t_2}(\C e(s),\EE \dot w(s))_2\,\de s.\label{eq:lowbou3}   
\end{align}
If we sum~\eqref{eq:lowbou1} and~\eqref{eq:lowbou2} and we use~\eqref{eq:lowbou3} we deduce that
\begin{align*}
&\mathcal Q(e(t_2)-e(t_1))+\tilde Q(\alpha(t_1),p(t_2)-p(t_1))+\|\nabla p(t_2)-\nabla p(t_1)\|_2^2+  \|\nabla \alpha(t_2)-\nabla \alpha(t_1)\|_2^2\\
&\le \int_{t_1}^{t_2}(\C (e(s)-e(t_1)),\EE \dot w(s))_2\,\de s+D(\alpha(t_1))-D(\alpha(t_2))+(\dot{d}(\alpha(t_1)),\alpha(t_2)-\alpha(t_1))_2\\
&\quad+\tilde Q(\alpha(t_1),p(t_2))-\tilde Q(\alpha(t_2),p(t_2))+(\dot{\B}(\alpha(t_1))(\alpha(t_2)-\alpha(t_1))p(t_1),p(t_1))_2.
\end{align*}
By~\eqref{eq:d} and~\eqref{eq:B1}, for all $\alpha_1,\alpha\in [0,1]$ and $\xi\in\M$ we have
\begin{align*}
|d(\alpha_1)-d(\alpha_2)-\dot{d}(\alpha_1)(\alpha_1-\alpha_2)| &\le \|\ddot{d}\|_\infty|\alpha_1-\alpha_2|^2,\\
|\B(\alpha_1)\xi:\xi-\B(\alpha_2)\xi:\xi-\dot{\B}(\alpha_1)(\alpha_1-\alpha_2)\xi:\xi| &\le \|\ddot{\B}\|_\infty|\alpha_1-\alpha_2|^2|\xi|^2.
\end{align*}
Therefore, thnkas to~\eqref{eq:C2} and the uniform estimates~\eqref{eq:unif_bound} (see also Remark~\ref{rem:unif_bound}), we have
\begin{align}
&\frac{\gamma_2}{2}\|e(t_2)-e(t_1)\|_2^2+\tilde Q(\alpha(t_1),p(t_2)-p(t_1))+\|\nabla p(t_2)-\nabla p(t_1)\|_2^2+  \|\nabla \alpha(t_2)-\nabla \alpha(t_1)\|_2^2\nonumber\\
&\le \int_{t_1}^{t_2}(\C (e(s)-e(t_1)),\EE \dot w(s))_2\,\de s+\|\ddot{d}\|_\infty\|\alpha(t_2)-\alpha(t_1)\|_2^2+\|\ddot{\B}\|_\infty\|\alpha(t_2)-\alpha(t_1)\|_4^2\|p(t_1)\|_4^2\nonumber\\
&\quad+\tilde Q(\alpha(t_1),p(t_2))-\tilde Q(\alpha(t_2),p(t_2))-\tilde Q(\alpha(t_1),p(t_1))+\tilde Q(\alpha(t_2),p(t_1)),\nonumber\\
&=\int_{t_1}^{t_2}(\C (e(s)-e(t_1)),\EE \dot w(s))_2\,\de s+\|\ddot{d}\|_\infty\|\alpha(t_2)-\alpha(t_1)\|_2^2+\|\ddot{\B}\|_\infty\|\alpha(t_2)-\alpha(t_1)\|_4^2\|p(t_1)\|_4^2\nonumber\\
&\quad+((\B(\alpha(t_1))-\B(\alpha(t_2))(p(t_2)-p(t_1)),p(t_2)+p(t_1))_2\nonumber\\
&\le C_1\|\alpha(t_2)-\alpha(t_1)\|_4^2+C_1\|\alpha(t_1)-\alpha(t_2)\|_4\|p(t_2)-p(t_1)\|_2+C_1\int_{t_1}^{t_2}\|e(s)-e(t_1)\|_2\|\EE \dot w(s)\|_2\,\de s\label{eq:part1}
\end{align}
for a constant $C_1>0$.

By Remark~\ref{rem:BCS}, the uniform estimates~\eqref{eq:unif_bound}, and~\eqref{eq:rH},~\eqref{eq:lowbou3}, there exists a constant $C_2>0$ such that
\begin{align*}
r_H\|p(t_2)-p(t_1)\|_1&\le \mathcal H(p(t_2)-p(t_1))\\
&\le \mathcal E(\alpha(t_1),e(t_1),p(t_1))-\mathcal E(\alpha(t_2),e(t_2),p(t_2))+\int_{t_1}^{t_2}(\C e(s),\EE \dot w(s))_2\,\de s\\
&\le \mathcal Q(e(t_1))-\mathcal Q(e(t_2))+D(\alpha(t_1))-D(\alpha(t_2))+\|\nabla \alpha(t_1)\|_2^2-\|\nabla \alpha(t_2)\|_2^2\\
&\quad+\tilde Q(\alpha(t_1),p(t_1))-\tilde Q(\alpha(t_1),p(t_2))+\tilde Q(\alpha(t_1),p(t_2))-\tilde Q(\alpha(t_2),p(t_2))\\
&\quad+\|\nabla p(t_1)\|_2^2-\|\nabla p(t_2)\|_2^2+\int_{t_1}^{t_2}(\C e(s),\EE \dot w(s))_2\,\de s\\
&\le C_2\left(\|e(t_2)-e(t_1)\|_2+\|\alpha(t_2)-\alpha(t_1)\|_4+\|\nabla\alpha(t_2)-\nabla\alpha(t_1)\|_2\right)\\
&\quad+C_2\left(\sqrt{\tilde Q(\alpha(t_1),p(t_2)-p(t_1))}+\|\nabla p(t_2)-\nabla p(t_1)\|_2+\int_{t_1}^{t_2}\|\EE \dot w(s)\|_2\,\de s\right).
\end{align*}
Therefore, if we use also~\eqref{eq:part1} we deduce the existence of a constant $C_3>0$ such that
\begin{align}
&r_H^2\|p(t_2)-p(t_1)\|_1^2\nonumber\\
&\le 3C_2^2\left(\|e(t_2)-e(t_1)\|_2^2+\|\alpha(t_2)-\alpha(t_1)\|_4^2+\|\nabla\alpha(t_2)-\nabla\alpha(t_1)\|_2^2\right)\nonumber\\
&\quad+3C_2^2\left(\tilde Q(\alpha(t_1),p(t_2)-p(t_1))+\|\nabla p(t_2)-\nabla p(t_1)\|_2^2+\left(\int_{t_1}^{t_2}\|\EE \dot w(s)\|_2\,\de s\right)^2\right)\nonumber\\
&\le C_3\|\alpha(t_2)-\alpha(t_1)\|_4^2+C_3\|\alpha(t_1)-\alpha(t_2)\|_4\|p(t_2)-p(t_1)\|_2\nonumber\\
&\quad+C_3\int_{t_1}^{t_2}\|e(s)-e(t_1)\|_2\|\EE \dot w(s)\|_2\,\de s+C_3\left(\int_{t_1}^{t_2}\|\EE \dot w(s)\|_2\,\de s\right)^2.\label{eq:part2}
\end{align}
By summing~\eqref{eq:part1} and~\eqref{eq:part2} and using also Corollary~\ref{coro:H1norm}, we can find a constant $C_4>0$ such that
\begin{align*}
&\|e(t_2)-e(t_1)\|_2^2+\|p(t_2)-p(t_1)\|_{H^1}^2+\|\alpha(t_2)-\alpha(t_1)\|_{H^1}^2\\
&\le C_4\|\alpha(t_2)-\alpha(t_1)\|_4^2+C_4\|\alpha(t_1)-\alpha(t_2)\|_4\|p(t_2)-p(t_1)\|_2\\
&\quad+C_4\int_{t_1}^{t_2}\|e(s)-e(t_1)\|_2\|\EE \dot w(s)\|_2\,\de s+C_4\left(\int_{t_1}^{t_2}\|\EE \dot w(s)\|_2\,\de s\right)^2.
\end{align*}
Firstly, we use Young Inequality
$$\|\alpha(t_1)-\alpha(t_2)\|_4\|p(t_2)-p(t_1)\|_2\le C_\epsilon\|\alpha(t_1)-\alpha(t_2)\|_4^2+\epsilon\|p(t_2)-p(t_1)\|_{H^1}^2\quad\text{for all $\epsilon>0$},$$
and then Lemma~\ref{lem:Cdelta} with $\theta=4$ to derive that
\begin{align*}
&\|e(t_2)-e(t_1)\|_2^2+\|p(t_2)-p(t_1)\|_{H^1}^2+\|\alpha(t_2)-\alpha(t_1)\|_{H^1}^2\\
&\le C_5\|\alpha(t_2)-\alpha(t_1)\|_1^2+C_5\int_{t_1}^{t_2}\|e(s)-e(t_1)\|_2\|\EE \dot w(s)\|_2\,\de s+C_5\left(\int_{t_1}^{t_2}\|\EE \dot w(s)\|_2\,\de s\right)^2
\end{align*}
for a constant $C_5>0$. In particular, there exist a constant $C>0$ such that
\begin{align}
&\left(\|e(t_2)-e(t_1)\|_2+\|p(t_2)-p(t_1)\|_{H^1}+\|\alpha(t_2)-\alpha(t_1)\|_{H^1}\right)^2\nonumber\\
&\le C^2\|\alpha(t_2)-\alpha(t_1)\|_1^2+C\int_{t_1}^{t_2}\|e(s)-e(t_1)\|_2\|\EE \dot w(s)\|_2\,\de s+\left(C\int_{t_1}^{t_2}\|\EE \dot w(s)\|_2\,\de s\right)^2.\label{eq:continuity}
\end{align}
Let us define the functions
\begin{align*}
&f(t)\coloneqq\|e(t)-e(t_1)\|_2+\|p(t)-p(t_1)\|_{H^1}+\|\alpha(t)-\alpha(t_1)\|_{H^1}\quad\text{for all $t\in[t_1,T]$},\\
&g(t)\coloneqq C\|\alpha(t)-\alpha(t_1)\|_1\quad\text{for all $t\in[t_1,T]$},\qquad h(t)\coloneqq C\|\EE \dot w(t)\|_2\quad\text{for a.e.\ $t\in[t_1,T]$}.
\end{align*}
Therefore, by~\eqref{eq:continuity} we deduce
\begin{align*}
f(t)^2\le g(t)^2+\int_{t_1}^tf(s)h(s)\,\de s+\left(\int_{t_1}^th(s)\,\de s\right)\quad\text{for all $t\in[t_1,T]$}.
\end{align*}
We can apply Lemma~\ref{lem:gronw} to deduce~\eqref{eq:cont1}. Finally, by Korn Inequality we have
$$\|u(t_2)-u(t_1)\|_2\le C_6\|e(t_2)-e(t_1)\|_2+C_6\|p(t_2)-p(t_1)\|_2+C_6\|w(t_2)-w(t_1)\|_{H^1}$$
for a constant $C_6>0$. Hence, by~\eqref{eq:cont1} we derive~\eqref{eq:cont2}.

\end{proof}

\end{document}
